\section{引言}

随着对更快、更高效计算的不懈追求，光子计算平台正受到越来越多的关注，这源于光学系统相较于电子系统固有的优势：大规模并行性、超低延迟和极低的能量耗散\cite{shastri2021photonics,prucnal2016recent,wetzstein2020inference}。在这些平台中，超表面——亚波长结构的纳米光子界面——已成为紧凑型光学计算的核心元件，因为它可以将波前调控、空间频率滤波、偏振控制和光谱选择性压缩到单层结构化表面中\cite{yu2014,koenderink2015,genevet2017,kamali2018review,chen2022review}。在模拟图像处理领域，超表面能够实现对光波前的直接操控，完成边缘检测、图像锐化和特征提取等关键任务，而无需光电转换\cite{silva2014,zhu2017,guo2018laplace,kwon2018,cordaro2019}。这些能力为自主导航、生物医学诊断和增强/虚拟现实系统等新兴应用奠定了基础\cite{wang2020metasurface,zhou2022flat}。与传统的笨重$4f$光学系统（依赖于成对透镜和空间滤波器）相比，基于超表面的处理器在尺寸、重量和集成潜力方面具有显著优势\cite{goodman2017,bykov2018optica}。特别是介质超表面，能够在保持亚波长分辨率的同时实现高透过效率，使其成为片上光学计算的理想候选方案\cite{decker2015high,west2014all}。

超表面光学计算的物理机制是光学传递函数（OTF）的工程化设计，OTF描述了光学系统如何调制入射光场不同空间频率分量的振幅和相位\cite{goodman2017,wesemann2021review}。在角谱表示中，任何输入图像都可以分解为具有横向波矢$k_x$和$k_y$的平面波分量，它们与入射角$\theta$和方位角$\phi$的关系为
\[
k_x = k_0 \sin\theta \cos\phi, \qquad
k_y = k_0 \sin\theta \sin\phi, \qquad
k_0 = \frac{2\pi}{\lambda},
\]
其中$k_0$是自由空间波矢\cite{goodman2017}。对于相干、线性且空间不变的光学系统，输出场$U(x,y)$由OTF $H(k_x, k_y)$通过傅里叶变换关系决定：
\[
U(x,y) = \mathcal{F}^{-1}\left[H(k_x, k_y) \cdot \mathcal{F}[U_0(x,y)]\right].
\]
其中$U_0(x,y)$是入射场，$\mathcal{F}$表示傅里叶变换，$\mathcal{F}^{-1}$是逆傅里叶变换。通过精心设计OTF $H(k_x, k_y)$，各种数学运算可以直接在光学域实现。例如，一阶空间微分器要求OTF正比于$ik_x$，其在角谱中呈现奇对称性，能够实现相应方向的边缘检测\cite{silva2014,bykov2018optica}。二阶微分器要求OTF正比于$k_x^2$，对应于拉普拉斯运算，可各向同性地增强边缘\cite{zhu2017,zhou2025laplace}。更复杂的运算如四阶微分\cite{qiu2025highorder}和高斯高通滤波\cite{zhou2023tailoring}可以通过设计OTF对横向波矢的高阶多项式或指数依赖关系来实现。

尽管早期的超表面设计主要关注在固定波长和正入射条件下实现所需的光学功能，但OTF工程的全部潜力在于控制光响应在波长（$\lambda$）、入射角（$\theta$）和偏振（$p$）完整参数空间上的分布。在本文中，我们提出了角度分辨傅里叶算子超表面（AFOM）的概念，其中OTF被显式地视为这三个变量的函数：
\[
\mathrm{OTF} = H(k_x, k_y; \lambda, \theta, p).
\]
在此框架下，各种光学操作作为该多维响应空间中的特定目标几何构型而涌现。空间微分要求透射系数遵循特定的角度依赖关系：
\[
t(\theta) \propto \theta, \qquad
t(\theta) \propto \theta^2, \qquad
t(\theta) \propto \theta^4,
\]
分别对应一阶微分\cite{silva2014}、二阶微分\cite{zhu2017,kwon2018}和四阶微分\cite{qiu2025highorder}。这些操作利用精心设计超原子单元的非局域共振响应，在角谱上创建所需的多项式OTF形状。偏振复用操作利用OTF的张量特性，使不同偏振通道能够进行独立处理\cite{deng2024,cotrufo2023natcom}。通过设计各向异性超原子，可以将OTF设计为$H_{xx} \neq H_{yy}$，使正交偏振分量同时经历不同的数学运算。时空微分将OTF概念扩展到频域，其中超表面响应被设计为同时依赖于空间频率（$k_x$）和时间频率（$\omega$），从而能够对时变信号进行运算\cite{tan2022,cordaro2019}。卷积算子可以通过设计与空间域特定核函数相对应的OTF来实现，从而实现超越微分的更一般线性图像处理\cite{zhou2023tailoring}。AFOM范式将这些多样化的功能统一在一个共同的数学框架下：
\[
\mathbf{R}^\star = \mathbf{R}^\star(\lambda, \theta, p),
\]
其中每种操作都对应于波长-角度-偏振空间中的特定目标响应张量，而逆向设计问题就变成了寻找能够实现这些目标OTF几何构型的结构。

尽管OTF工程在概念上很优雅，但AFOM的逆向设计面临着制约实际实现的根本性挑战。从超表面结构到光学响应的映射是高维的、高度非线性的，且通过严格的电磁仿真评估计算成本极高\cite{molesky2018}。对于具有数千个活动像素的自由形态二值超表面，设计空间巨大，结构-响应关系表现出复杂的依赖关系。OTF敏感地依赖于几何参数（周期、填充因子、厚度、形状），以耦合的、不可分离的方式变化\cite{sell2017}。给定波长和角度下的共振响应通过Kramers-Kronig关系和因果性约束与整个光谱的行为相关联\cite{wesemann2021review}。非局域超表面表现出强烈的角色散，入射角的微小变化可以显著改变透射相位和振幅\cite{guo2018laplace}。用于偏振复用的各向异性结构可能遭受不期望的交叉偏振转换，降低操作保真度\cite{deng2024}。传统设计方法依赖于经验试错或在有限几何族（如简单光栅、孔阵列或棒状结构）内的参数扫描\cite{kwon2018,cordaro2019}。当针对复杂OTF形状、宽带操作或大数值孔径时，这些方法变得不切实际。此外，逆向问题本质上是多对一的：多个物理上不同的结构可能都在容差范围内满足相同的目标OTF，这表明确定性优化方法可能会错过更优的替代方案\cite{piggott2015,piggott2017}。

机器学习方法已成为超表面逆向设计的有前途工具，提供了从数据中学习复杂结构-响应映射的能力\cite{ma2021deep,liu2018training,so2019}。现有方法可以分为正向代理模型（预测给定几何参数的光学响应）\cite{ma2021deep,malkiel2018plasmonic}、直接逆向网络（将目标响应映射回几何参数）\cite{liu2018training,an2019deep}、生成模型（包括GANs\cite{liu2018,ma2021}、VAEs\cite{ma2019probabilistic}和扩散模型\cite{hen2026,tanriover2023}，学习有效设计的分布）以及物理信息方法（将电磁约束纳入学习框架）\cite{chen2021physics,huang2022meta}。然而，当应用于AFOM设计时，这些现有方法表现出关键局限性。大多数方法局限于简单、固定的结构模板，只有少量可调参数，无法处理具有数千个自由度的自由形态二值结构\cite{ma2021deep,malkiel2018plasmonic}。现有的逆向设计框架通常针对单一功能进行训练，缺乏处理不同目标OTF的灵活性\cite{an2019deep,jiang2021}。大多数先前的工作仅关注正入射或窄角度范围，忽略了高NA操作和现实成像系统所需的全角度分辨响应\cite{kwon2018,zhou2025laplace}。许多方法将目标简化为标量指标或稀疏光谱点，丢弃了完整波长-角度-偏振响应张量中的结构化信息\cite{liu2018training}。生成的设计可能违反制造要求，需要事后修改，从而降低性能\cite{phillips2021ensemble}。

本工作通过开发用于AFOM自由形态逆向设计的物理一致生成框架来解决现有方法的局限性。首先，我们构建了一个全面的制造感知自由形态二值超表面结构数据集，通过严格的RCWA模拟获得涵盖波长、入射角和偏振通道的OTF标签。数据生成流程在结构合成期间强制执行物理制造约束（最小特征尺寸、连通性），确保生成设计的实际可实现性。其次，我们提出了一个闭环设计工作流程，包括用于目标OTF规范的任务特定评分函数、具有自适应质量阈值的物理一致数据整理、条件扩散模型训练用于结构生成、具有物理引导推理的代理辅助候选生成、以及基于仿真验证的重排序进行最终选择。该流程确保报告的设备与验证的散射行为相关联，而非仅依赖神经网络置信度。第三，我们证明了多样化的光学操作——二阶微分、偏振无关操作、偏振复用处理、四阶微分、低通滤波和时空微分——可以在相同响应空间内统一为不同的目标几何构型。通过条件化适当的目标OTF张量$\mathbf{R}^\star(\lambda, \theta, p)$，同一训练框架可以处理所有这些任务。最后，我们实现了最先进的实验基准：二阶微分器具有多波长操作、高数值孔径（NA $\approx$ 0.4）、宽带性能和高透射效率（验证任务分数超过0.93）；偏振复用设计可对正交偏振进行独立微分运算；四阶和低通滤波器具有复杂OTF形状（验证分数0.993）；时空微分器可实现空间和时间信息的联合处理。

本文其余部分组织如下。第2节介绍统一响应公式和完整的物理一致设计框架。第3节报告跨基准任务的全面实验结果。第4节讨论扩展趋势、局限性和更广泛的影响。第5节给出结论。
